\cotefile{138}
\scanpage{17}
le sous-corps (ou mieux, sous-anneau) de $\overline{Q}$
qu'ils engendrent, ext. finie de $Q$. On voit de
suite qu'il est indépendant du choix de
$s$, et que c'est la plus petite sous-ext\unsure{ension}
de $\overline{Q}/Q$ telle que $\widehat{\mathscr X}$ puisse se descendre \inserted{en $\widehat X_K$} : $K$
de façon que les él. de $S$ \unsure{deviennent} rat/$K$.
De plus la descente en question est unique,
à isom unique près. Les \unsure{autom\textsuperscript{ismes}|automor\textsuperscript{phismes}} de
$(\widehat{\mathscr X},S)$ (ou encore de $\mathscr X$) font un groupe
fini ($\hookrightarrow \unsure{\mathfrak S|\mathscr S}_S$), et ce groupe opère \unsure{bel-|bt.}\unsure{\& bien|déjà}
sur $\widehat X_K$.
% archive sheet number: 16
